\documentclass[11pt]{article}
\usepackage{amsmath}
\usepackage{amssymb}
\usepackage{fontspec}
\usepackage{unicode-math}
\setmainfont{DejaVu Serif}
\setsansfont{DejaVu Sans}
\setmonofont{DejaVu Sans Mono}
\setmathfont{DejaVu Math TeX Gyre}
\usepackage[a4paper,margin=2.5cm]{geometry}
\usepackage{booktabs}
\usepackage{longtable}
\usepackage{array}
\usepackage{enumitem}
\usepackage{xcolor}
\usepackage[colorlinks=true,linkcolor=blue!50!black,urlcolor=blue!50!black,citecolor=blue!50!black]{hyperref}
\usepackage{parskip}
\usepackage{fancyhdr}
\pagestyle{fancy}
\fancyhf{}
\renewcommand{\headrulewidth}{0.4pt}
\fancyhead[C]{\footnotesize ARob -- UAVs / Aeronaves Robotizadas}
\fancyfoot[C]{\thepage}
\setlength{\headheight}{26pt}
\sloppy
\emergencystretch=2em
\title{UAVs Labs - Knowledge Base (2025/26 edition)}
\author{Course notes -- 2025/26 labs}
\date{}
\begin{document}
\maketitle
\tableofcontents
\vspace{1em}

Entry point into a deep technical analysis of the three lab assignments (Parrot AR.Drone) and their DevKit material.
Detailed, equation-by-equation notes for each lab live in \texttt{notes/} (linked below).
This file gives the lab overview, how the three labs chain together, and how they connect to the theory in \href{../lectures-22-23/COURSE_NOTES.pdf}{\texttt{../lectures-22-23/COURSE\_\allowbreak{}NOTES.pdf}}.

\section{Course/lab identity}

\begin{itemize}
\item \textbf{Course}: Unmanned Aerial Vehicles / Aeronaves Robotizadas, MSc Aerospace Engineering (MEAer), Instituto Superior Técnico, Lisboa - 2025/26 First Semester.
\item \textbf{Platform}: Parrot AR.Drone 2.0 (indoor hull), controlled over its own Wi-Fi access point (\texttt{192.168.1.1}), via MATLAB/Simulink + the "ARDrone Simulink Development Kit" (David Escobar Sanabria \& Pieter J. Mosterman, U. Minnesota / MathWorks, v1.0, Sept. 2013 - unmodified vendor devkit reused as-is across all three labs and, per the lecture notes' course logistics, across editions).
\item \textbf{Deliverable format (identical for all 3 labs)}: single-column PDF report, \textbf{max 15 pages}, plus the \texttt{.m} scripts, submitted via Fenix, using the shared cover page \texttt{25\_\allowbreak{}26\_\allowbreak{}UAVs\_\allowbreak{}Lab\_\allowbreak{}report\_\allowbreak{}cover\_\allowbreak{}page.docx} (Instituto/course/semester header + a 3-row group-member table + instructor field - only the "Laboratory ?" number changes between submissions).
\item \textbf{Grading weight} (from \href{../lectures-22-23/COURSE_NOTES.pdf}{\texttt{../lectures-22-23/COURSE\_\allowbreak{}NOTES.pdf}}): $L = 0.2 L1 + 0.4 L2 + 0.4 L3$, contributing 60\% of the final grade $F = 0.4 E + 0.6 L$.
\item \textbf{Question tagging convention (all labs)}: \textbf{(T)} = theoretical, solved before the lab session; \textbf{(L)} = laboratory, requires simulation/experimental data collected during the session.
\end{itemize}

\section{Lab map and detailed notes}

\begin{center}
\small
\setlength{\tabcolsep}{3pt}
\begin{longtable}{|p{1.1cm}|p{4.9cm}|p{3.1cm}|p{3.3cm}|p{2.6cm}|}
\hline
Lab & Focus (per course syllabus) & Handout & Devkit entry point & Detailed notes \\
\hline
\endhead
Lab 1 & Modelling and identification of the drone & \texttt{25\_\allowbreak{}26\_\allowbreak{}AR\_\allowbreak{}Lab1.pdf} (Sept. 2025) & \texttt{DevKit\_\allowbreak{}Lab1/.../start\_\allowbreak{}here.m} & \href{notes/01-lab1.pdf}{\texttt{notes/01-lab1.md}} \\
\hline
Lab 2 & Estimation of motion variables & \texttt{25\_\allowbreak{}26\_\allowbreak{}UAVs\_\allowbreak{}Lab2.pdf} (Oct. 2025) & \texttt{DevKit\_\allowbreak{}Lab2/.../start\_\allowbreak{}here\_\allowbreak{}NAV.m} & \href{notes/02-lab2.pdf}{\texttt{notes/02-lab2.md}} \\
\hline
Lab 3 & Motion control of the drone & \texttt{25\_\allowbreak{}26\_\allowbreak{}UAVs\_\allowbreak{}Lab3.pdf} (Nov. 2025) & \texttt{DevKit\_\allowbreak{}Lab3/.../start\_\allowbreak{}here\_\allowbreak{}CTRL.m} & \href{notes/03-lab3.pdf}{\texttt{notes/03-lab3.md}} \\
\hline
\end{longtable}
\end{center}

Supporting reference paper (required reading for Lab 2 Section 5): P. Batista, C. Silvestre, P. Oliveira, "Partial attitude and rate gyro bias estimation: observability analysis, filter design, and performance evaluation," \textit{International Journal of Control}, 84(5), 2011 - summarized in \href{notes/02-lab2.pdf}{\texttt{notes/02-lab2.md}}.

\section{How the three labs chain together}

The devkit itself grows incrementally across the three labs, and each lab's report is built directly on the previous one's results:

\begin{enumerate}
\item \textbf{Lab 1 (Modelling)} starts from the plain devkit (\texttt{start\_\allowbreak{}here.m}; models: baseline, hover/position, waypoint-tracking). Its theory section (§2) has students \textit{derive by hand} the mixer matrix, rotor-thrust allocation, hover linearization, and per-axis transfer functions that the course's Ch1/Ch2 lecture material sets up - notably, \textbf{Lab 1's Q2.2 is the concrete, worked version of the "× configuration" exercise that the Ch2 lecture slides explicitly left unsolved} (see the gotchas list in \href{../lectures-22-23/COURSE_NOTES.pdf}{\texttt{../lectures-22-23/COURSE\_\allowbreak{}NOTES.pdf}}). Its lab sections then \textit{experimentally identify} the height and pitch closed-loop dynamics (step response + frequency response) and compare against the six pre-identified linear models (\texttt{ssRoll}, \texttt{ssPitch}, \texttt{ssYaw}, \texttt{ssH}, etc.) that ship with the devkit and are reused, unmodified, in every later lab.
\item \textbf{Lab 2 (Estimation)} swaps in an upgraded devkit (\texttt{start\_\allowbreak{}here\_\allowbreak{}NAV.m}, adds a raw-navdata decoder \texttt{decode\_\allowbreak{}32bit.m} at 200 Hz, and an offline \textbf{Replay} mode) that exposes raw accelerometer/gyro samples that Lab 1's devkit did not. It explicitly starts from the same accelerometer-inclinometer baseline documented in the sensors lecture notes, builds up steady-state Kalman filters and bias-augmented complementary filters axis-by-axis (pitch, then roll), and culminates in implementing the full 6-state Kalman filter from the Batista/Silvestre/Oliveira (2011) paper - the rigorous, GAS generalization of the course's own "rate-gyro-only is unobservable, add a vector measurement to fix it" worked example.
\item \textbf{Lab 3 (Control)} swaps in a third devkit variant (\texttt{start\_\allowbreak{}here\_\allowbreak{}CTRL.m}, adds a \textbf{Trajectory Tracking} mode with two Simulink blocks, \texttt{desired\_\allowbreak{}trajectory} and \texttt{tt\_\allowbreak{}controller}, deliberately left empty for students to implement). Its theory section (§2) re-derives an LQR trajectory-tracking law plus an adaptive disturbance-rejection term (textbook instance of the Ch7 adaptive-control pattern, generalized from a scalar unknown parameter to a vector disturbance), then §3 has students implement it against straight-line and circular reference trajectories, finishing with a line-of-sight (LOS) path-following law - the concrete worked version of the "lookahead guidance" alternative that Ch8's lecture notes mention but never derive.
\end{enumerate}

Across all three labs, the same in-code warning appears verbatim in the Wi-Fi control scripts: \textbf{the AR.Drone's onboard position estimate is inaccurate}, because it integrates a noisy optical-flow-derived velocity estimate (the drone has no GPS - confirmed in \href{../lectures-22-23/notes/03-sensors-estimation.pdf}{\texttt{../lectures-22-23/notes/03-\allowbreak{}sensors-\allowbreak{}estimation.pdf}}). This single limitation is the throughline that motivates the whole lab sequence: Lab 1 identifies how bad the raw dynamics/estimates are, Lab 2 builds better estimators to compensate, and Lab 3's controller performance is explicitly expected to depend on how good an estimator Lab 2 produced.

\section{Consolidated gotchas / things to double-check before relying on this material}

\begin{itemize}
\item The devkit's \texttt{guide.pdf} and all \texttt{.m} scripts are unmodified since \textbf{2013} (MATLAB R2013b, Real-Time Windows Target - since renamed/deprecated to "Simulink Real-Time" in later MATLAB releases) despite being reused in the 2025/26 edition; worth confirming with the instructor whether a newer toolchain is actually expected.
\item \texttt{Devkit\_\allowbreak{}Lab1}'s and \texttt{Devkit\_\allowbreak{}Lab3}'s \texttt{guide.pdf} are \textbf{byte-identical} (verified via md5sum) and only document the original 3 examples (baseline, hover, waypoint-tracking) - it was never updated to document Lab 3's Trajectory Tracking mode; the Lab 3 handout itself is the only documentation for \texttt{start\_\allowbreak{}here\_\allowbreak{}CTRL.m}, \texttt{ARDroneTTSim.slx}/\texttt{ARDroneTT.slx}, and the \texttt{desired\_\allowbreak{}trajectory}/\texttt{tt\_\allowbreak{}controller} blocks.
\item No lab's \texttt{.m} scripts contain numeric controller/filter gains ($K$, $Q$, $R$, $kw$, LOS $\Delta$/$V$, Kalman $Q$/$R$, etc.) - by design, these are left for students to derive and tune; there is no "reference answer" in the source material for any of them. Lab 3 also reuses the symbol \texttt{kw} for two \textit{different} gains (an adaptation-law gain in §2.5 vs. a vertical proportional gain in §3.1) - worth flagging explicitly if writing a report.
\item The six pre-identified linear models (\texttt{ssRoll}, \texttt{ssRoll2V}, \texttt{ssPitch}, \texttt{ssPitch2U}, \texttt{ssYaw}, \texttt{ssH}) are shared, unmodified \texttt{.mat} files across all three devkits; their numeric contents were not inspected (binary) by any of the analysis forks - load them in MATLAB if a report needs the literal identified plant coefficients.
\item All \texttt{.slx} Simulink models are binary (zipped XML) and were not parsed for internal block-diagram/gain content in any lab - every claim about what a model does is inferred from its filename, the referencing \texttt{.m} scripts, and the handout text, not from opening the model itself.
\item \texttt{lib/getWaypoints.m}'s header comment (all 3 labs) describes a 4-field waypoint but the code actually builds a 5-column matrix - a minor, harmless vendor doc/code mismatch, not course-introduced.
\end{itemize}

See each \texttt{notes/0N-labN.md} file for full task-by-task detail, exact equations, and lab-specific open questions.

\end{document}
